\chapter[Sermon 70. The Fourth and Fifth Angels Shed Their Vials.]{The Fourth and Fifth Angels Shed Their Vials.}
\label{sermon:070}
\markright{Sermon 70. The Fourth and Fifth Angels Shed Their Vials.}

And the fourth angel poured out his vial upon the sun, and it was given to him to afflict men with heat of fire. And men raged with great heat, and blasphemed the name of God who had power over the plagues, and they did not repent to give him glory. And the fifth angel poured out his vial upon the seat of the beast, and his kingdom grew dark, and they gnawed their tongues for sorrow, and blasphemed the God of heaven because of the pain and anguish of their sores, and they did not repent of their deeds.

The faithful do not regard their afflictions, sent by God’s just judgment, as the punishments of sinners, but as trials of their faith; although they acknowledge that they are justly afflicted for their sins. Yet, they commend God’s grace, which transforms the punishments of sinners into trials of faith. To the ungodly, punishments are plagues, which they cannot bear patiently, nor glorify God, but rather blaspheme him, and believe they suffer undeservedly. Therefore, God’s plagues are most grievous to the ungodly, although far more cruel things are prepared for them, namely, everlasting damnation in the next world. Therefore, the plagues of this world inflicted upon the ungodly are like certain preparations and introductions to more grievous torments.

\index{Antichrist}\index{Apocalypse (Book of)}The fourth Angel pours out his vial on the sun, and to him was given power to plague people with heat or fire. Many interpret this plague allegorically, understanding by the sun Christ enlivening the consciences of the faithful, and the latter to be darkened in the minds of people, choosing rather the darkness of Antichrist than the light of Christ. Therefore, consciences, erring and seduced with error, burn with various lusts and despairs, by which they are driven at length to various blasphemies. While I do not entirely reject this interpretation, in my judgment the sense will be clearer if we understand the fourth plague to be heat and great drought, a barrenness of the earth, a scarcity of grain, finally an intolerable thirst afflicting both people and beasts, and also breeding and intensifying hot diseases. For so we have read in the threatenings of the law: I will give a heaven of brass and an earth of iron. In the time of Elijah, for despising and rejecting the word of the Lord, God plagued Israel with a severe drought, as you may see in the third book of Kings, chapters 17 and 18. Jeremiah also describes a similar drought and heat in chapter 14. Again, the Lord defended Israel with a pillar of a cloud by day and a pillar of fire by night. Moreover, we have heard previously in the Apocalypse: the sun shall not shine upon them, neither any heat. And justly is this world plagued with burning heat, as it offends grievously, burns with various lusts, and also by wicked proclamations prohibits the spreading and refreshing of God’s word.

For this is the effect of the plague. And he says to me, burned with great heat. At first, he says, being inflamed with an exceedingly great heat, they were even raging mad. For we read in histories that those afflicted with too much heat have felt grievous displeasures and torments both of body and mind. Then he adds what follows from the former passage: the impatience of the heat provoked them to blaspheme God, and even him who has power over these plagues; namely, that because he has full power to do so, he will not deliver them so vexed with burning heat. Conversely, the children of Israel in their tents being stung by serpents, inflaming the whole body with the sting, did repent and did not blaspheme God. But coming to Moses, they said, “We have sinned, for we have spoken against the Lord and against you. Pray the Lord that he will take away from us these serpents.” Therefore, those who, through impatience, murmur against the judgments of God blaspheme the name of the Lord, and will not acknowledge themselves to be rightly and justly punished, seeking pardon. Finally, it is added that they did not repent so that they might give glory to God, and so on. For the Lord plagues us to the end that, being afflicted, we should repent and give God the glory, confessing, as I said before, that we are justly punished, and ought with weeping and wailing to turn to the Lord striking us. But these, like Pharaoh, do not acknowledge their sin, nor pray to God, nor are they amended, but many times overcome themselves in malice. Hereof we learn the difference between the godly and ungodly, and how both use themselves in afflictions. For they give glory to God and amend their life; these do not give God the glory, but become worse than themselves. To give God the glory is to give place to God, not to resist, but to acknowledge their sin and God’s righteousness; and not only this, but also the mercy of God and clemency toward the penitent, and the same to require humbly.

\index{Papacy}\index{Rome}The fifth angel pours his cup upon the seat of the beast. It is more evident than it needs to be proven by testimonies that a seat or throne is used for a kingdom, as John himself immediately uses “seat” to mean a kingdom. Also, in times past, the masters, or rather ministers of churches taught while seated, and had their stoles and chairs in holy assemblies. The saying in the gospel is well known: “In the chair of Moses sit the scribes and Pharisees,” etc. It is known that in ancient times there were seats of patriarchs—Jerusalem, Antioch, Rome, Alexandria, Constantinople, and others—and that these are called Apostolic seats, because the Apostles taught there. And so the Apostolic seat is used to refer to the Apostolic doctrine itself. That seat erected and established at Rome by the Apostles and Apostolic men, the beast—that is, the Pope—has subverted, and in its place has erected the seat of pestilence, which he dares nevertheless to call the seat of Christ and the seat of Peter. Christ has no more any seat on earth, except that he dwells in the hearts of the faithful church. Otherwise, the true seat of Christ is the right hand of the Father. The true seat of Peter is heaven itself. Rome is no longer his seat, for the Apostolic doctrine and patriarchal chair have been destroyed and trodden under foot, and in their stead an earthly empire or kingdom has been set up by the Pope. Indeed, he overthrows the Apostolic seats by force of arms. Now, therefore, God, having compassion on his people, pours out his wrath and plague on the see of Rome, illuminating men with the light of the Gospel, so that they may know and see the wickedness and abomination of the Roman see. This is a wonderful benefit to those who are enlightened, and a great grief and torment to the Roman people. For the effect of the plague follows: and his kingdom was made dark. This plague corresponds to the ninth of Egypt. For just as thick darkness plagued the Egyptians, bright light rejoiced the Israelites, so the Papists will be tormented by shameful errors, and it will grieve them to have their errors detected and their glory obscured; the faithful will rejoice in the light of Christ. For now the majesty of the seat, and of him who sits therein, begins—and has already begun—to be obscured. What was once called a holy seat is now, by the godly and learned, called wicked Rome, the whore of Babylon, the mother of all fornications, the den of thieves, Sodom, Egypt, the red harlot by reason of the purple senate of Cardinals, who wear red and purple. It is commonly said and truly, the nearer Rome, the further from Christ. They rightly call the Cardinals, bishops, and spiritual fathers the family and limbs of Antichrist, men deceived and deceivers, most corrupt with simony and filthy lust. Therefore, the kingdom of the beast—so he explains the seat—was made dark. Furthermore, it is added how the worshippers of the seat of the beast behave themselves. First, for pain and sorrow, indignation, wrath, and envy, they gnaw or bite their tongues, which is the gesture of angry men, and impotently angry, that is, burning with furious rage. This is a phrase of speech signifying how they will rage with great fury against the truth that has been opened, which they would have utterly hidden and suppressed. Again, they blaspheme the Lord of heaven and maker of all, both because he afflicts them with sores and sundry plagues, and also because he casts a darkness upon their kingdom. For even now, the Romans call the preachers of the gospel deceivers and heretics, and the very doctrine of the gospel, heresy. But this reproach redounds to him who is the author of that doctrine. Finally, they do not repent of their doings, of their simony, of their crafty juggling, sacrileges, idolatry, and all ungodliness. And the apostle says that evil men and deceivers will grow worse and worse, deceiving and being deceived. Therefore, it is no marvel that you see the Papists at this day, with a stiff neck, proceeding obstinately in their errors. But the greatest plague is to be forsaken by God, and stubbornly to maintain their errors and ungodliness, and therein to persevere. The Lord deliver us from evil. Amen.
